\section{引言：从无条件到条件生成}

本讲是 KAIST CS492D（Diffusion Models and Their Applications）的第六讲，由 Minhyuk Sung 教授讲授。前几讲已经介绍了 DDPM 和 DDIM 等扩散模型的基本原理，学生已具备实现基础扩散模型的能力。本讲将"换挡"，讨论如何在实际应用中引入\textbf{条件信息}（class label、文本描述、图像等），以及如何\textbf{高效微调}大规模预训练扩散模型。

\begin{knowledgebox}{前置知识回顾}
在 DDPM/DDIM 框架中，噪声预测网络 $\epsilon_\theta(x_t, t)$ 的输出与 score function 之间存在确定性关系：
$$
\nabla_{x_t} \log p(x_t) \;\propto\; -\frac{1}{\sqrt{1-\bar\alpha_t}}\,\epsilon_\theta(x_t, t)
$$
其中 $\bar\alpha_t$ 为噪声调度的累积乘积。这一关系是所有条件生成方法的数学基础——通过修改 score function 来引导生成轨迹。
\end{knowledgebox}

本讲覆盖以下核心主题：
\begin{enumerate}
    \item \textbf{Classifier Guidance}：利用外部分类器引导条件生成
    \item \textbf{Classifier-Free Guidance (CFG)}：无需分类器的条件引导，工业标配
    \item \textbf{DDIM Inversion}：确定性反演与图像编辑
    \item \textbf{Latent Diffusion Models}：在潜在空间运行扩散以降低计算开销
    \item \textbf{ControlNet}：高效微调实现图像条件控制
    \item \textbf{LoRA}：低秩适应用于扩散模型个性化
\end{enumerate}

\newpage
